% ------------------------------------------------------------
% Page 105
% ------------------------------------------------------------

L’Ecole fut fermée en 1940-41. Elle fut réouverte en 1942. Mon successeur, M. B. de
Villefort, utilisa la salle de classe mais son épouse ne souhaita pas occuper la sacristie
et le ménage trouva à se loger dans le hameau.
Pour clore ce chapitre, ennuyeux pour d’autres que moi-même j’indique que le
pasteur et sa jeune femme sont bien venus célébrer le culte en 1938 devant une salle
comble.

J’ajoute que l’école fut définitivement fermée en 1944, que le temple fut désaffecté et
acquis, peu après par les frères Rousset qui, après aménagement, en firent leur maison
d’habitation où je leur rendis visite quelques années plus tard.

Le hameau, dès la fin de la guerre, a connu le destin de tous les villages de
campagne ; les Rousset, eux-mêmes se retirèrent à Pourcarès, hameau moins éloigné
de Meyrueis. La maison forestière de Vaballe fut abandonnée mais quelques jolis
chalets de pêcheurs, occupés à la belle saison, se sont dressés au bord de la route qui
domine la jolie rivière à truites qu’est restée la Brèze.

Je suis bien conscient que mes réflexions et hypothèses m’ont emmené bien loin de
mon projet initial : « Vous faire part des problèmes auxquels nous fûmes confrontés,
ma femme et moi durant notre séjour de 9 mois dans notre « temple-école ». Mais aux
souvenirs qu’on voudrait évoquer s’attachent d’autres souvenirs qu’il est difficile de
dissocier et d’éluder ;

\begin{center}
\textbf{La vie quotidienne : problèmes à résoudre :}
\end{center}

\textbf{Le ravitaillement}

Facilement résolu. Un épicier ambulant « le Caïffa » livrait à domicile, tous les
samedis, tout ce qui était essentiel pour la semaine : pain, vin, légumes verts ou secs,
fruits et tout un arsenal de conserves sans oublier les bidons de pétrole dont nous
faisions une grande consommation durant les mois d’hiver.
Pour d’autres petits achats je descendais quand même assez souvent à Meyrueis et, au
moins toutes les fins de mois, pour retirer à la perception mes 1030 f. d’émoluments
mensuels. Mon épouse, peu douée pour le vélo, descendit à pied au chef-lieu mais
revint épuisée par les 25 km de l’aller et retour. Elle ne renouvela pas l’expérience et
ne connut de Meyrueis que nos passages rapides lors de nos départs et retours de
vacances.

\textbf{Le chauffage}

A neuf cent mètres et plus d’altitude nous subissions un froid intense dans notre petit
deux pièces et d’autant plus que cuisine et chambrette étaient dépourvues de volets
protecteurs. Mon vieux pardessus d’Ecole Normale étendu devant « le fenestrou »
pour nous protéger de la lumière, était souvent le matin, selon l’expression consacrée,
aussi « raide que la justice ».
